\documentclass[11pt,a4paper]{article}

\usepackage[T1]{fontenc}
\usepackage[utf8]{inputenc}
\usepackage{lmodern}
\usepackage{microtype}
\usepackage[margin=2.2cm]{geometry}
\usepackage{amsmath,amssymb}
\usepackage{booktabs,tabularx,longtable,array}
\usepackage[table]{xcolor}
\usepackage{enumitem}
\usepackage{fancyhdr}
\usepackage{underscore}
\usepackage[hidelinks]{hyperref}

\definecolor{accent}{HTML}{2C4A7C}
\definecolor{soft}{HTML}{E9EFF7}
\definecolor{warnsoft}{HTML}{FBF1DF}
\definecolor{warnline}{HTML}{8A5A11}

\hypersetup{colorlinks=true,linkcolor=accent,urlcolor=accent,citecolor=accent}

\setlist{itemsep=2pt,topsep=4pt,parsep=0pt}
\renewcommand{\arraystretch}{1.18}
\setlength{\LTpre}{6pt}
\setlength{\LTpost}{6pt}

\pagestyle{fancy}
\fancyhf{}
\fancyhead[L]{\small\color{accent}Rerun protocol --- STL decomposition study}
\fancyhead[R]{\small\color{accent}v1.7 --- approved (v1.4) + change log}
\fancyfoot[C]{\small\thepage}
\renewcommand{\headrulewidth}{0.4pt}

\newcommand{\code}[1]{\texttt{#1}}
\newcommand{\RN}{\mathrm{RelNaive}^{(10)}}
\newcommand{\dstl}{\Delta_{\mathrm{STL}}}
\newcommand{\dft}{\Delta_{\mathrm{FT}}}

\setlength{\fboxsep}{8pt}
\setlength{\headheight}{14pt}
\newsavebox{\notebox}
\newenvironment{keybox}%
  {\par\medskip\noindent\begin{lrbox}{\notebox}\begin{minipage}{\dimexpr\linewidth-2\fboxsep-2\fboxrule\relax}}%
  {\end{minipage}\end{lrbox}\fcolorbox{accent}{soft}{\usebox{\notebox}}\par\medskip}
\newenvironment{decisionbox}%
  {\par\medskip\noindent\begin{lrbox}{\notebox}\begin{minipage}{\dimexpr\linewidth-2\fboxsep-2\fboxrule\relax}}%
  {\end{minipage}\end{lrbox}\fcolorbox{warnline}{warnsoft}{\usebox{\notebox}}\par\medskip}

\title{\color{accent}\textbf{Rerun Protocol}\\[4pt]
\large When Does Seasonal Decomposition Help?\\
Pipeline decisions, measurement and analysis plan, implementation,\\ testing, pilot and full runs}
\author{Prepared for D.\ Costa, the supervisor, and the partner university team}
\date{Version 1.7 --- 3 October 2026 --- v1.4 approved by the supervisor; later amendments in Appendix~\ref{app:changelog}}

\begin{document}
\maketitle
\thispagestyle{fancy}

\begin{keybox}
\textbf{What this document is.} The single reference for the rerun of the study. It records every
design decision and the reason for it, fixes the complete list of quantities to be measured and the
analyses to be run on them \emph{before} any data is produced, and lays out how the pipeline is
implemented from the existing code, how it is tested, how the pilot is run and judged, and how the
full runs are executed. After supervisor sign-off it is committed and tagged in the pipeline
repository and frozen; later changes go only into the change log (Appendix~\ref{app:changelog})
with a date and a reason.

\smallskip
\textbf{Bottom line.} No existing forecast file is reused as a result. The existing code is reused
after the changes in Section~\ref{sec:impl}. One pilot, then one set of runs from one frozen
manifest and one frozen configuration file, in strict mode, on four GPU machines. The analysis code
is written and tested on pilot data before the full run starts.
\end{keybox}

\begin{decisionbox}
\textbf{Status of the decisions.} Every decision in Section~\ref{sec:decisions} was taken by
D.\ Costa on 28 September 2026 in eight rounds of structured questions. The rejected alternatives
and their reasons are recorded in Appendix~\ref{app:alternatives}. The five technical
specifications marked ``Protocol'' in the register (the eligibility rule derived from the length
constraints, decomposition details, search spaces, the statistical cache and provenance fields)
were not asked separately and are approved with the protocol at sign-off. The supervisor approves this protocol before implementation and the pilot's
go / no-go before the full run.
\end{decisionbox}

\tableofcontents
\newpage

%=====================================================================
\section{Purpose, scope and verdict on existing results}\label{sec:purpose}
%=====================================================================

\subsection{Why a rerun}

A review of the forecasting pipeline and of every run it produced (28 September 2026) found defects
in the data the paper was built on and in the design of the experiment itself. The full list is in
Appendix~\ref{app:defects}; Table~\ref{tab:defects-short} summarises those that force a rerun.

\begin{table}[h]
\centering\small
\caption{Defects that force a rerun. IDs are used throughout this document.}\label{tab:defects-short}
\begin{tabularx}{\linewidth}{@{}l X X@{}}
\toprule
ID & Defect & Consequence \\
\midrule
C1 & The paper's global baselines and all statistical STL results came from a non-strict run whose
forecasts are a hard-coded blend of seasonal naive, drift and mean, jittered by $\pm0.5\%$ per task.
& Original Tables 3--8 do not describe any trained model. \\
C2 & No global model is tuned. The ``Auto'' models are fixed configurations; tuning code writes empty
placeholder files. The paper claims Optuna tuning over 100 trials. & Methods section is inaccurate;
statistical models (self-selecting) are compared with untuned learners. \\
C3 & Features are computed on the series minus only the last horizon, so they see the test targets
of windows 0 and 1. & Bucket assignment, and every ``which series'' result, uses test data. \\
C4 & The ML seasonal lag and the neural input size are set by the \emph{shortest} series in each
training pool; only 33\% of quarterly cohort fits had the seasonal lag. & Both experiments are
confounded by whether the model saw seasonal information. \\
M1 & The two trained bodies of results use different series, and no trained global STL-AC run
exists. & Specialisation cannot be measured on a paired basis. \\
F1 & Features silently fall back to fixed values (0.0 or 1.0) on short series; 1.7--2.9\% of the
old sample. & Those series pile into the extreme buckets, e.g.\ 51 series in the Low
evolving-seasonality bucket. \\
\bottomrule
\end{tabularx}
\end{table}

\subsection{Verdict: what is reused and what is not}

\begin{itemize}
  \item \textbf{Not reused as results:} any forecast or metric file from the 23 May, 25 May or
  28--30 May runs. The 23 May run is not model output (C1). The other two are genuine but carry C3
  and C4, and fixing C4 changes the forecasts themselves; mixing old and new files would recreate
  the unpaired design (M1).
  \item \textbf{Reused as a check:} the 25 May strict run's statistical forecasts. Statistical fits
  do not depend on the training pool, so the new code must reproduce them row for row on the same
  series (integration test I2).
  \item \textbf{Reused as code:} data loading, STL transform, rolling-origin split, RelNaive metric,
  the six feature definitions, checkpointing, scheduler and vast.ai scripts --- after the changes in
  Section~\ref{sec:impl}, in the new folder \code{pipeline/} of the project repository.
  \item \textbf{Reused as text:} introduction, related work, the 64 verified references, problem
  formulation and the structure of the methods and results sections.
\end{itemize}

\subsection{Principles}\label{sec:principles}

Every decision in Section~\ref{sec:decisions} follows from one of six principles.

\begin{enumerate}[label=\textbf{P\arabic*}]
  \item \textbf{One manifest, one configuration, one code version.} Every forecast in the study is
  traceable to a single frozen manifest, a single frozen configuration file and a single git commit.
  \item \textbf{Nothing from a test period is used before evaluation.} Features, sampling, buckets,
  tuning and scaling use only observations before the first test window.
  \item \textbf{A model's configuration never depends on who else is in the pool.} Lags, input size
  and all hyperparameters come from the frozen configuration; series too short for them are excluded
  before sampling.
  \item \textbf{Every family receives comparable care.} Statistical models select their own
  specification; global models are tuned with a comparable, documented budget.
  \item \textbf{No silent failure path.} There is no fallback forecaster in the code. A failure is
  recorded as a failure and counted in the results.
  \item \textbf{The analysis refuses dirty inputs.} Automated gates check provenance, hashes,
  completeness and pairing before any number is computed.
\end{enumerate}

%=====================================================================
\section{Decision register}\label{sec:decisions}
%=====================================================================

``Team'' marks decisions taken by D.\ Costa in the structured questions of 28 September 2026;
``Protocol'' marks technical specifications approved with the protocol at sign-off. Rejected
alternatives are in Appendix~\ref{app:alternatives}; details in
Sections~\ref{sec:data}--\ref{sec:analysis}.

{\small
\begin{longtable}{@{}p{0.6cm} p{2.1cm} p{6.3cm} p{4.3cm} p{1.5cm}@{}}
\toprule
ID & Topic & Decision & Rationale & By \\
\midrule
\endfirsthead
\toprule
ID & Topic & Decision & Rationale & By \\
\midrule
\endhead
\bottomrule
\endfoot
\multicolumn{5}{@{}l}{\emph{Scope and data}}\\
D1 & Reuse & Full rerun. Legacy statistical forecasts used only as a regression check (I2). & Section~\ref{sec:purpose}. & Team \\
D2 & Data & M3 and M4, monthly ($m=12$, $h=18$) and quarterly ($m=4$, $h=8$) only. One frozen copy, stored with SHA-256 checksums; every machine loads the same bytes. M3 from the original competition file \code{M3C.xls} with official ids, cross-checked against the Monash and Mcomp copies; M4 from the official competition files (amended in v1.3 and v1.4). & No new frequencies or loaders to validate; a changed mirror cannot change the data; the M3 data exactly as published. & Team \\
D3 & Eligibility & A series is eligible if its history before the first test window has at least $L = 3m+h$ points: 54 monthly, 20 quarterly. & Every STL fit sees $\geq 2$ seasons; tuning data holds $\geq 3m$ points; evolving seasonality gets $\geq 2$ rolling windows. Fixes C4, M4. & Protocol \\
D4 & Features & Definitions unchanged. Computed once, on the pre-window-0 history, for every eligible series. A series with any non-``ok'' quality flag is ineligible. & Fixes C3 and F1; buckets fixed across windows; the follow-up gets features for the whole pool. & Team \\
D5 & Feature guards & Option A: quarterly minimum lengths lowered to 20 (nonlinearity from 30, ARCH from 24); monthly guards ($\leq 48$) unchanged. Sensitivity analysis A7. & Keeps 600 per source and three windows (Section~\ref{sec:eligibility}). & Team \\
D6 & Sampling & Per-feature stratified quantile sampling, 30 strata, 600 series per source, one seed per feature, candidate pool = all eligible series. & The pool must exceed the quota for stratification to exist (M2). & Team \\
D7 & Buckets & Terciles per feature $\times$ frequency within the feature's sample. Degenerate if any populated bucket holds $< 20\%$: run with populated buckets, excluded from trend tests. Nonlinearity keeps its floored definition and is expected to be flagged. & Makes the collapse explicit; keeps the paper's failure-mode contribution. & Team \\
D8 & Validation & Rolling origin, $W = 3$ windows, step $h$; every model refitted every window. & Standard protocol; all families treated alike. & Team \\
\multicolumn{5}{@{}l}{\emph{Models and tuning}}\\
D9 & Decomposition & Robust STL (statsmodels 0.14.6), fitted per window on training data only. Strategies Direct, STL-SN, STL-AC. STL failure is a failure; no fallback decomposition. & Fixes F2. & Protocol \\
D10 & Models & The 12 models: ARIMA, SARIMA ($D=1$), ETS; Ridge, LinearRegression, RandomForest, XGBoost; NLinear, NHITS, LSTM; TFT, PatchTST. & Design and text stay aligned. & Team \\
D11 & Statistical & AutoARIMA with approximate stepwise search, AutoARIMA with $D=1$, AutoETS (statsforecast 2.0.3); specification chosen per series by information criteria. & As in the valid May run; regression test I2 stays exact. & Team \\
D12 & Global tuning & Library Auto classes with Optuna (TPE); 20 trials per study; the slowest family may drop to 15 only if the pilot projects an overrun, recorded before any real run. Objective: mean validation MASE. & Fixes C2; MASE is scale-free and native. & Team \\
D13 & Tuning set & Fixed draw of 600 per source from the eligible pool; pre-window-0 data only; validation = one block, the last $h$ points before the first cutoff. One study per model $\times$ frequency $\times$ target (raw, $T{+}R$, and $T$, $S$, $R$ separately): 90 studies. Best configuration frozen to a hashed file. & Tuning never sees test data; each STL-AC component model is tuned for its own target. & Team \\
D14 & Search spaces & Table~\ref{tab:spaces}; fixed a priori; identical for every pool; include scaling and training length. & P3; fixes C4, M5. & Protocol \\
D15 & Training scopes & Cohort: one global model per feature sample $\times$ frequency, M3 and M4 pooled (1,200 series). Tercile: one per bucket ($\approx 400$). Same frozen configuration, same series, same windows; temporal hold-out on the same series. & Fully paired design; corrects the ``held-out series'' wording (M3). & Team \\
D25 & ML inputs & Lags $1,\dots,m$; target transform tuned (none, local scaling, first difference, seasonal difference); no calendar or rolling features. & Same information as a seasonal autoregression; STL's contribution measured cleanly. & Team \\
\multicolumn{5}{@{}l}{\emph{Execution}}\\
D16 & Failures & No fallback code. Statistical per-series failures recorded as failed rows. A global task is retried once; a second failure is fixed in a committed change and that task alone is rerun with the same manifest and configurations, logged. & P5; fixes C1, m2. & Team \\
D17 & Seeds & One fixed seed per fit. Seed check: evolving seasonality, cohort scope, the 7 stochastic global models, 2 extra seeds (A15). & Quantifies seed sensitivity at modest cost. & Team \\
D18 & Statistical cache & Statistical forecasts computed once per series $\times$ window $\times$ strategy $\times$ model, reused across feature samples. & Pool-independent; saves compute; identical values. & Protocol \\
D19 & Provenance & Every row stores code commit, environment-lock hash, data checksum, manifest hash, configuration hash, backend and version, and --- for STL-AC --- all three component provenances. & Fixes m1, F3. & Protocol \\
D20 & Infrastructure & Four GPU machines, sharded by family $\times$ frequency; outputs and logs synced continuously to a cloud bucket via rclone, with checksums. & Survives instance loss; merge refuses mismatched hashes. & Team \\
D22 & Code \& environment & The pipeline in a new top-level folder \code{pipeline/} of the project repository, developed on branch \code{rerun-v2} from an unmodified baseline commit (amended in v1.2); the May library versions pinned in a committed lockfile (Section~\ref{sec:env}). & One history for paper and code; every fix is a visible diff; regression test valid. & Team \\
\multicolumn{5}{@{}l}{\emph{Analysis}}\\
D21 & Metrics & Primary: $\RN$. Secondary: unclipped RelNaive, MASE, sMAPE, POCID. & Continuity with the paper. & Team \\
D23 & Inference & Series-level Wilcoxon signed-rank with Holm correction per predefined table; multiple comparison matrices; Friedman + Nemenyi critical-difference rankings for cohort (36 configurations) and tercile (27) scopes, overall and per frequency; bootstrap CIs; ordered-trend tests. & Section~\ref{sec:analysis}. & Team \\
D26 & Practical threshold & An effect is a \emph{gain} or \emph{loss} only if Holm-significant and $|$median per-series $\Delta| \geq 0.01$; significant but smaller effects are \emph{negligible}. & With $\sim$10$^6$ rows almost everything is significant. & Team \\
D27 & Aggregation & Descriptives at instance level; inference at distinct-series level; families weighted equally when pooled. & Paper's convention; no series or family counted twice. & Team \\
D28 & Selector & No feature-based selector in this paper; stated as a limitation and left to the follow-up. & Scope. & Team \\
\multicolumn{5}{@{}l}{\emph{Process}}\\
D24 & Freeze & Protocol committed and tagged in the pipeline repository before stage R1; later changes only via the change log. & Prevents result-driven changes. & Team \\
D29 & Release & Code, manifest and frozen result bundle released with the paper. & Every table and figure regenerable by reviewers. & Team \\
D30 & Sign-off & The supervisor approves the protocol before implementation and the pilot's go / no-go; the partner university is informed. & Accountability at both gates. & Team \\
\end{longtable}
}

%=====================================================================
\section{Data, eligibility, features and sampling}\label{sec:data}
%=====================================================================

\subsection{Data}

M3 and M4 monthly and quarterly series are downloaded once, validated, stored in the cloud bucket
together with SHA-256 checksums, and loaded from that frozen copy on every machine (D2). The checksum
is written into the manifest and every result row. Loading verifies it and refuses a mismatch
(test U18).

\emph{Sources (D2, amended in v1.3 and v1.4).} M3 is taken from the original competition file
\code{M3C.xls} of the International Institute of Forecasters (M3 competition page), which holds the
training and test values of every series, the official ids (\code{N0646}, \dots) and start periods.
The site refuses scripted downloads, so the file is supplied to the freeze, which refuses it unless its
SHA-256 equals the value pinned in the code and stores it with the raw files. Two further copies are
compared with it value by value, and the freeze stops on any difference that is not documented: the
Monash Time Series Forecasting Archive copy on Zenodo, which the earlier runs used, is identical at full
precision (204,566 of 204,566 values); Mcomp~2.7 (CRAN; R.~J.~Hyndman) differs in one value, series
N2786 at $t=84$ (a test value), where it has $1200$ and the original $-1200$, without a note in its
change log. The earlier runs read the Monash file through \code{datasetsforecast}, whose reader
rounded 44,034 values to single precision (relative error ${\le}\,5.7\times10^{-8}$; defect m7); the
freeze reads it at full precision. Series are named by their official M3 ids; the rows of
\code{M3C.xls} are checked to be in official-id order, so the positional ids of the earlier runs
(\code{M1}, \dots) map one-to-one onto official ids, as Mcomp's field \code{st} confirms. The start
labels are kept as given, except that N1071's ``quarter 9 of 1984'' is carried into 1986-Q1; the 29
monthly series without a start date keep their positions as labels. M4 is taken from the official
competition files (\code{Mcompetitions/M4-methods}, \code{Dataset/Train} and \code{Dataset/Test}) and read
by \code{read_m4_official}\footnote{\code{src/data/freeze_data.py}; \code{M4.load} of
\code{datasetsforecast} melts the whole wide table and was stopped by the out-of-memory killer on a
machine with 7.6\,GiB.}, which returns exactly what \code{M4.load} returns (checked against
\code{M4.load} itself in the tests and, for M4 Quarterly, on the real files) and in addition refuses a
training row with a gap, a test row without exactly $h$ finite values, and ids that are duplicated or
not matched one-to-one.

\emph{Validation and layout.} Both sources are complete series: the competition's test period is part
of every series, and the rolling-origin windows of Section~\ref{sec:forecasting} are cut from them.
The freeze stops unless every source has its published number of series (M3: 1,428 monthly, 756
quarterly; M4: 48,000 and 24,000) and its shortest series equals the published minimum training
length plus $h$ (M3: 48+18 and 16+8; M4: 42+18 and 16+8), which confirms that the series are
complete. Each source is stored as one file with columns \code{unique_id}, \code{source_dataset},
\code{frequency}, position \code{t} $=1,\dots,n$, \code{y} and the source's time label (kept for
traceability only). The manifest (\mbox{\code{pipeline/data_manifest/m3m4_v1.json}}) records the SHA-256 of
every raw download and every frozen file, a content hash that does not depend on the file format,
the results of the M3 cross-checks, and the code commit and package versions that produced the copy.
The pipeline uses only the order of observations: the loader rebuilds one regular calendar from
\code{t}, identical for M3 and M4, and has no fallback of any kind (defect F4,
Appendix~\ref{app:defects}).

\subsection{Eligibility and the minimum-length constraint}\label{sec:eligibility}

With $W=3$ windows of step $h$, the history available before the first test window is $n - 3h$.
Table~\ref{tab:minlen} lists every requirement on that history. The eligibility rule $L = 3m+h$
satisfies all modelling requirements. The only conflicts were the quarterly guards of nonlinearity
and ARCH, resolved by D5.

\begin{table}[h]
\centering\small
\caption{Minimum pre-window-0 length required by each component.}\label{tab:minlen}
\begin{tabular}{@{}lrrl@{}}
\toprule
Requirement & Monthly & Quarterly & Source \\
\midrule
STL with $\geq 2$ seasons in the tuning data (pre-window-0 minus $h$) & 42 & 16 & D3 \\
ML lags $\leq m$ plus seasonal differencing at the tuning origin & 42 & 16 & D25 \\
Evolving seasonality: $\geq 2$ rolling windows (window + step) & 42 & 18 & feature code \\
Non-normality test & 20 & 20 & feature code \\
Structural breaks & 48 & 20 & feature code \\
ARCH test (quarterly: original 24, lowered by D5) & 48 & 20 & feature code \\
Nonlinearity (quarterly: original 30, lowered by D5) & 48 & 20 & feature code \\
\midrule
\textbf{Eligibility rule $L = 3m+h$} & \textbf{54} & \textbf{20} & D3 \\
\bottomrule
\end{tabular}
\end{table}

Table~\ref{tab:eligible} gives the eligible pool under the adopted option and the alternatives that
were considered, estimated from the 25 May run's cutoffs; exact counts are produced by stage R1.

\begin{table}[h]
\centering\small
\caption{Estimated eligible series. A (adopted): three windows, quarterly guards lowered to 20.
B: two quarterly windows, guards of 24. C: three windows, original guards (30). D: three windows,
guards of 24. Monthly figures are the same under every option, because the monthly rule (54)
already exceeds every monthly guard (48).}\label{tab:eligible}
\begin{tabular}{@{}lrrrrr@{}}
\toprule
Source & Pool & A (adopted) & B & C & D \\
\midrule
M3 Monthly   & 1,428  & 1,040  & 1,040  & 1,040  & 1,040 \\
M3 Quarterly & 756    & \textbf{634} & 666 & 116 & 433 \\
M4 Monthly   & 48,000 & 35,000 & 35,000 & 35,000 & 35,000 \\
M4 Quarterly & 24,000 & 22,500 & 23,400 & 20,100 & 20,700 \\
\bottomrule
\end{tabular}
\end{table}

Under option A, nonlinearity and ARCH are computed on 20--29 points for the shortest quarterly
series and are noisier there; analysis A7 repeats the feature results on the series that meet the
original guards. M3 Quarterly is close to a census (600 of about 634 eligible series) and M3
Monthly supplies 600 of about 1,040, so per-feature stratification is substantive only for M4. The
paper states this.

\subsection{Features}

The six feature definitions of the paper are unchanged, with three changes to how they are computed
(D4).
\begin{enumerate}
  \item Each feature is computed once, on $y_{1:n-3h}$, the history before the first test window.
  Features and buckets are the same for all three windows.
  \item Any series with a non-``ok'' quality flag on any feature is removed from the eligible pool,
  so no fallback value ever enters a sample (fixes F1).
  \item Features are computed for \emph{every} eligible series, about 59,000, so the follow-up
  meta-learning paper has all six features for every eligible series.
\end{enumerate}

\subsection{Sampling}\label{sec:sampling}

For each feature $k$, frequency and source independently:
\begin{enumerate}
  \item Rank the eligible series by $\phi_k$ (ties broken by series id).
  \item Cut the ranking into 30 equal-count strata.
  \item Draw $600/30 = 20$ series per stratum without replacement, seed $= s_0 + k$.
  \item Record the stratum, rank and quantile of every selected series in the manifest.
\end{enumerate}
Each feature sample holds 600 series per source: 1,200 per frequency and 2,400 in total. Across
the six features, about 8,500 distinct series are expected. M4 samples are nearly disjoint; M3
samples overlap heavily. Sampling quality is checked by gate G11.

\subsection{Buckets}

Within each feature sample and frequency, the 1,200 series are split at the empirical 1/3 and 2/3
quantiles into Low, Medium and High, about 400 each. If any populated bucket holds less than 20\%
of the sample, the feature is flagged degenerate for that frequency (gate G12). It is then run with
its populated buckets only and excluded from the tercile-trend tests (H4). The nonlinearity
feature keeps its floored definition (D7) and is expected to trigger this guard, as documented in
the paper's Section 4.5.

\subsection{The manifest}

Stage R1 writes one manifest bundle, which is hashed and frozen. It contains:
\begin{itemize}
  \item the data checksum;
  \item the eligibility table for the full pool (length, flags, eligible yes/no);
  \item feature values for all eligible series;
  \item the six feature samples with strata and seeds;
  \item bucket assignments;
  \item the tuning set;
  \item the cutoffs for all windows.
\end{itemize}
Nothing downstream re-samples, re-buckets or recomputes cutoffs.

%=====================================================================
\section{Forecasting setup}\label{sec:forecasting}
%=====================================================================

\subsection{Validation and decomposition}

For series $i$ and window $w \in \{0,1,2\}$, the training set ends at $n_i - (3-w)h$ and the test
set is the next $h$ points. Every model --- statistical and global --- is refitted for every
window (D8). STL is fitted on each window's training data only. Its three strategies are:
\begin{itemize}
  \item \textbf{Direct:} the model forecasts the raw series.
  \item \textbf{STL-SN:} a model forecasts $T+R$, and the last seasonal cycle of $S$ is repeated and
  added back.
  \item \textbf{STL-AC:} one separately trained model per component, and the forecasts are summed.
\end{itemize}

\subsection{Models and classes}

\begin{table}[h]
\centering\small
\caption{Model classes used in each phase.}\label{tab:models}
\begin{tabularx}{\linewidth}{@{}l l l X l@{}}
\toprule
Family & Model & Tuning phase & Evaluation phase & Type \\
\midrule
Statistical & ARIMA & --- (IC selection) & \code{AutoARIMA(season_length=m)}, approximate & local \\
Statistical & SARIMA & --- & \code{AutoARIMA(season_length=m, D=1)}, approximate & local \\
Statistical & ETS & --- & \code{AutoETS(season_length=m)} & local \\
ML & Ridge / LinReg / RF / XGBoost & \code{AutoMLForecast} + model spaces & \code{MLForecast} with frozen config & L / L / NL / NL \\
Neural & NLinear, NHITS, LSTM & \code{AutoNLinear}, \code{AutoNHITS}, \code{AutoLSTM} & plain classes, frozen config & L / NL / NL \\
Transformer & TFT, PatchTST & \code{AutoTFT}, \code{AutoPatchTST} & plain classes, frozen config & NL / NL \\
\bottomrule
\end{tabularx}
\end{table}

The Auto classes are used only in the tuning phase. The evaluation phase instantiates the plain
classes from the frozen configuration file, so every evaluation forecast is fully determined by that
file. The Auto class names and signatures are verified in the pinned environment on day 1 (test I5),
because the libraries are not installed on the development machine.

\subsection{Tuning protocol}\label{sec:tuning}

\begin{itemize}
  \item \textbf{Data:} the tuning set (D13), truncated at each series' first cutoff. The last $h$
  points form a single validation block; the rest is training data. No tuning step sees a test
  period (gate G7).
  \item \textbf{Studies:} one per global model $\times$ frequency $\times$ target, i.e.\
  $9 \times 2 \times 5 = 90$. The five targets are the raw series (Direct), $T+R$ (STL-SN), and $T$,
  $S$, $R$ separately (STL-AC). 50 studies run on GPU (neural and transformer) and 40 on CPU (ML).
  \item \textbf{Search:} Optuna TPE sampler, 20 trials, fixed sampler seed. If the pilot projects
  an overrun, the slowest family may drop to 15 trials; the change is recorded before any real run
  (D12). Studies are stored in SQLite and archived.
  \item \textbf{Objective:} mean validation MASE with seasonality $m$. NeuralForecast is configured
  with \code{valid_loss=MASE(seasonality=m)}; AutoMLForecast with an equivalent loss function.
  \item \textbf{Early stopping:} used during tuning only. In evaluation, the number of training steps
  is part of the frozen configuration, since no validation data exists after a test cutoff.
  \item \textbf{Freezing:} the best trial of each study is written to \code{configs_frozen.json},
  whose SHA-256 hash is stored in every result row. Tercile-scope models use exactly the same
  configuration as cohort-scope models (D15).
\end{itemize}

\begin{table}[h]
\centering\small
\caption{Search spaces (D14). They are identical for every pool. The pilot may only \emph{narrow}
a space for feasibility, and any such change goes into the change log.}\label{tab:spaces}
\begin{tabularx}{\linewidth}{@{}l X@{}}
\toprule
Scope & Space \\
\midrule
All neural and transformer & \code{input_size}: monthly $\{18, 36\}$, quarterly $\{8, 16\}$; \code{scaler_type} $\in$ \{identity, standard, robust\}; learning rate log-uniform $[10^{-4}, 10^{-2}]$; \code{max_steps} $\in \{500, 1000, 2000\}$; batch size $\in \{32, 64\}$; \code{start_padding_enabled}=True (fixed). \\
NLinear & common parameters only \\
NHITS & \code{n_pool_kernel_size} $\in$ \{[2,2,1], [1,1,1]\}; \code{n_freq_downsample} $\in$ \{[4,2,1], [1,1,1]\}; MLP units $\in$ \{2$\times$64, 2$\times$256\} \\
LSTM & encoder hidden $\in \{32, 64, 128\}$; encoder layers $\in \{1, 2\}$; decoder hidden $\in \{32, 64\}$ \\
TFT & hidden size $\in \{32, 64, 128\}$; heads $\in \{2, 4\}$; dropout $\in \{0, 0.1, 0.2\}$ \\
PatchTST & hidden size $\in \{32, 64, 128\}$; heads $\in \{4, 8\}$; encoder layers $\in \{1, 2, 3\}$; patch length $\in \{4, 8\}$, capped at \code{input_size}; stride = patch length / 2 \\
All ML (D25) & lags $1,\dots,m$ (fixed); target transform $\in$ \{none, local standard scaling, first difference, seasonal difference\}; no calendar or rolling features \\
Ridge & $\alpha$ log-uniform $[10^{-3}, 10^{2}]$ \\
LinearRegression & target transform only \\
RandomForest & trees $\in \{100, 300\}$; max depth $\in$ \{none, 10, 20\}; min samples per leaf $\in \{1, 5, 20\}$; max features $\in$ \{1.0, 0.5, sqrt\} \\
XGBoost & trees $[100, 600]$; learning rate log-uniform $[0.01, 0.3]$; max depth $[3, 9]$; subsample $[0.6, 1]$; column sample $[0.6, 1]$ \\
\bottomrule
\end{tabularx}
\end{table}

\subsection{Training scopes and pairing}

For each feature sample, frequency, strategy, model and window, the cohort model is trained on all
1,200 series of that frequency --- M3 and M4 pooled --- and the three tercile models on about 400
each. All four use the same frozen configuration and are evaluated on the same series and windows.
Every tercile forecast therefore has exactly one cohort twin, and gate G6 enforces this. $\dft$
isolates the effect of the narrower training pool.

Statistical models appear only in the cohort scope. Their forecasts do not depend on the pool, so
they are computed once per series, window, strategy and model and reused across feature samples
(D18).

\subsection{Failure semantics and determinism}

\begin{itemize}
  \item \textbf{No fallback forecaster.} \code{baselines.py}, the component stub and the
  per-task jitter are deleted. The result object rejects any source other than \code{trained_*}
  (test U9).
  \item \textbf{Statistical models:} a per-series exception produces a \emph{failed row} with the
  exception class and message, and no forecast. The task continues.
  \item \textbf{Global models:} an exception fails the whole task, which is retried once
  automatically. If it fails again, the cause is fixed in a committed code change and only that
  task is rerun, with the same manifest and configurations; the change log records it (D16).
  \item \textbf{STL-AC:} if any of the three component models fails, the row fails. All three
  component provenances are stored.
  \item \textbf{Seeds:} one fixed seed per fit, with deterministic backend flags where available.
  Forecast hashes are stored so that any rerun can be compared exactly. Seed sensitivity is
  measured by the seed check (D17, A15).
\end{itemize}

\subsection{Code, environment and data}\label{sec:env}

The pipeline lives in the new top-level folder \code{pipeline/} of the project repository
(\code{github.com/DiogoCosta18/TimeSeriesForecasting}) and is developed on branch \code{rerun-v2},
starting from an unmodified baseline commit; the paper and the code therefore share one history, and
the results record the commit that produced them (D22). The environment is pinned to the versions of the valid May runs:
Python 3.11.10, torch 2.5.1+cu121, mlforecast 1.0.31, neuralforecast 3.1.8, statsforecast 2.0.3,
optuna 4.8.0, statsmodels 0.14.6, scikit-learn 1.8.0, numpy 2.4.6, pandas 2.3.3, scipy 1.15.3,
ruptures 1.1.10 and datasetsforecast 1.0.1, plus rdata 0.11.2 for the Mcomp cross-check (added in v1.3); \code{M3C.xls} is read with xlrd 2.0.2. A lockfile is committed and its hash is stored in every
row. The same image is used on every machine, and all machines load the frozen data copy (D2).

%=====================================================================
\section{Measurement and analysis plan}\label{sec:analysis}
%=====================================================================

This section is frozen before the full run (D24). The analysis code implements it and is tested on
pilot data first (I3).

\subsection{What is recorded}

\begin{itemize}
  \item \textbf{Per forecast row} (task $\times$ series $\times$ window): forecast and
  seasonal-naive vectors; $\RN$ and unclipped RelNaive; MASE, sMAPE, MAE, POCID; the keys (feature
  sample, frequency, strategy, family, model, scope, bucket, series, source, window, cutoff, seed);
  training and inference time; status (trained / failed) with the failure reason; and the full
  provenance of D19.
  \item \textbf{Per series:} length, eligibility, all six features with flags, sample membership,
  stratum, bucket per feature.
  \item \textbf{Per tuning study:} every trial's parameters, validation MASE and duration; the best
  configuration; the SQLite study file.
  \item \textbf{Per run:} task-level logs, heartbeat, resource usage, gate reports.
\end{itemize}

\subsection{Derived quantities and aggregation rules}

\begin{itemize}
  \item \textbf{Instance:} (feature sample, series, window, model, strategy, scope, seed).
  \item $\dstl = \RN_{\text{Direct}} - \RN_{\text{STL}}$, paired on everything except the
  strategy. $\dft = \RN_{\text{cohort}} - \RN_{\text{tercile}}$, paired on everything except the
  scope. Positive values mean the treatment helps.
  \item \textbf{Descriptives} (mean of capped values, median, win rate = share of instances with
  $\Delta > 0$, tie share) are computed at the \emph{instance} level. This keeps the convention of
  the current paper (D27).
  \item \textbf{Inference} is at the \emph{distinct-series} level. A series' value is its mean over
  the instances in the reporting group --- windows, models and, where a series appears in several
  feature samples, those samples. No series is counted twice.
  \item \textbf{Family pooling:} when families are pooled, each family is weighted equally (mean of
  family means), so the four-model ML family does not outweigh the two transformers (D27).
  \item \textbf{Seed check:} only the main seed enters H1--H7; the extra seeds are used by A15 only.
  \item \textbf{Failures:} pairwise deletion. Every table reports the number of excluded instances.
\end{itemize}

\subsection{Practical significance}\label{sec:practical}

With about $2.8 \times 10^6$ forecast rows, almost every difference is statistically significant.
Every confirmatory contrast is therefore labelled by the rule of D26:
\begin{itemize}
  \item \textbf{Gain / loss:} Holm-significant and $|\text{median per-series } \Delta| \geq 0.01$,
  i.e.\ at least 1\% of the seasonal-naive error, in the direction of the sign.
  \item \textbf{Negligible:} Holm-significant but $|\text{median per-series } \Delta| < 0.01$.
  \item \textbf{No evidence:} not significant after Holm correction.
\end{itemize}
For H4 the contrast is the difference in median per-series $\dstl$ between the High and Low
terciles; for H6 it is the median of the per-series paired difference. The paper uses these three
labels verbatim; ``gain'' is never used for a negligible effect.

\subsection{Confirmatory hypotheses}

All tests are two-sided at $\alpha = 0.05$, with Holm correction within the stated family of tests.
Effect sizes are reported with every test: the median per-series $\Delta$ with a 95\% bootstrap
confidence interval (2,000 resamples of series) and the matched-pairs rank-biserial correlation.

{\small
\begin{longtable}{@{}p{0.6cm} p{5.1cm} p{5.1cm} p{4.0cm}@{}}
\toprule
ID & Question & Test and unit & Correction family \\
\midrule
\endfirsthead
\toprule
ID & Question & Test and unit & Correction family \\
\midrule
\endhead
\bottomrule
\endfoot
H1 & Does STL change accuracy, per family and strategy (cohort scope)? & Wilcoxon signed-rank on per-series $\dstl$ & Holm: 4 families $\times$ 2 strategies = 8 \\
H2 & Does the STL effect differ between families? & Paired Wilcoxon on per-series $\dstl$ between families (same series) & Holm: 6 pairs $\times$ 2 strategies = 12 \\
H3 & Which strategy is best within each family? & Paired Wilcoxon on $\RN$, Direct vs STL-SN vs STL-AC; shown as multiple comparison matrices & Holm: 3 pairs $\times$ 4 families = 12 \\
H4 & Does the STL effect vary with each feature? & (a) Jonckheere--Terpstra trend test of per-series $\dstl$ across Low $<$ Medium $<$ High; (b) Spearman $\rho$ between feature value and per-series $\dstl$, bootstrap CI. Family-balanced pooling; per family as secondary. & Holm: 6 features $\times$ 2 strategies = 12 per test type \\
H5 & Does bucket specialisation change accuracy, per family and strategy? & Wilcoxon on per-series $\dft$ & Holm: 3 families $\times$ 3 strategies = 9 \\
H6 & Do linear and non-linear models respond differently to specialisation? & Paired Wilcoxon on per-series [mean $\dft$ of non-linear models $-$ mean $\dft$ of linear models] & Holm: 3 strategies = 3 \\
H7 & Do configurations differ in overall rank? & Friedman test on per-series mean $\RN$ with distinct series as blocks, then Nemenyi post-hoc and critical-difference diagrams: cohort scope (36 configurations = 12 models $\times$ 3 strategies) and tercile scope (27 = 9 $\times$ 3), each overall and per frequency --- six diagrams. & Nemenyi's own family-wise control within each diagram \\
\end{longtable}
}

For H7, the critical difference shrinks with the number of blocks, so with thousands of series
most pairs will be separated in rank. A rank difference is read as practically meaningful only
where the corresponding paired contrast also meets the D26 threshold.

\subsection{Secondary and robustness analyses}

\begin{description}[leftmargin=1.1cm,style=sameline,font=\normalfont\bfseries]
  \item[A1] Model standing: $\RN$ by family $\times$ model $\times$ strategy $\times$ scope; top
  configurations; distributions.
  \item[A2] Frequency split: H1, H4 and H5 repeated for monthly and quarterly separately.
  \item[A3] Source split: M3 vs M4.
  \item[A4] Window consistency: sign and size of $\dstl$ and $\dft$ per window.
  \item[A5] Metric sensitivity: H1 and H5 on MASE, sMAPE and unclipped median RelNaive.
  \item[A6] Cap sensitivity: means at caps 5, 10 and 20, and 5\% trimmed means.
  \item[A7] Guard sensitivity (D5): H4 for nonlinearity and ARCH restricted to quarterly series with
  $\geq 30$ and $\geq 24$ pre-window-0 points respectively.
  \item[A8] Failure accounting by family, model, strategy and scope.
  \item[A9] Compute cost: training and inference time by family and strategy.
  \item[A10] Tuning report: frozen configurations, trial curves, best vs first trial.
  \item[A11] Sampling quality: Kolmogorov--Smirnov distance between sample and eligible pool per
  feature and source.
  \item[A12] Bucket support and degenerate-feature report.
  \item[A13] Per-model tables for $\dstl$ and $\dft$.
  \item[A14] Correlations between the six features on the eligible pool.
  \item[A15] Seed sensitivity (D17): for evolving seasonality in the cohort scope, the spread of
  per-series $\RN$ across three seeds for the seven stochastic global models, and whether any H1 or
  H3 label (gain / loss / negligible / no evidence) changes when another seed is used.
\end{description}

\subsection{Tables and figures, mapped to the paper}

\begin{table}[h]
\centering\small
\caption{Outputs produced by the analysis package. All are generated by committed scripts from the
frozen result bundle.}\label{tab:outputs}
\begin{tabularx}{\linewidth}{@{}l X l@{}}
\toprule
Paper section & Output & Analysis \\
\midrule
\S2 Data & Pool / eligible / sampled counts; sampling illustration; KS table & D3, A11 \\
\S2 Models & Model table; search spaces; frozen configurations (appendix) & D10--D14, A10 \\
\S2 Features & Feature distributions; bucket sizes with support check & A12, A14 \\
\S3.1 & Top configurations per scope; $\RN$ distributions by family $\times$ strategy; cohort-scope critical-difference diagrams & A1, H7 \\
\S3.2 & $\dstl$ by family $\times$ strategy with CIs and D26 labels; bar chart; MCM per family & H1--H3 \\
\S3.3 & $\dstl$ by feature $\times$ tercile (per strategy); trend and Spearman tables; family heatmap & H4, A7 \\
\S4 & $\dft$ overall, by bucket, family, model, strategy; linear vs non-linear; significance counts; tercile-scope critical-difference diagrams & H5--H7, A13 \\
Appendix & Frequency, source, window, metric, cap and seed robustness; failures; compute cost & A2--A9, A15 \\
\bottomrule
\end{tabularx}
\end{table}

%=====================================================================
\section{Implementation plan}\label{sec:impl}
%=====================================================================

\subsection{From scratch or from the current code}

\begin{table}[h]
\centering\small
\caption{Rewrite versus refactor.}\label{tab:scratch}
\begin{tabularx}{\linewidth}{@{}l X X@{}}
\toprule
 & From scratch & Refactor the current code (adopted) \\
\midrule
For & No hidden legacy path survives; cleaner architecture. & STL, split, metric, features, checkpointing and vast.ai scripts are already correct and tested; feature code \emph{is} the paper's definition. \\
Against & Four days is not enough to rebuild and re-test loading, features, checkpointing and orchestration; new bugs; features must be re-validated against the old definitions. & Every unsafe path must be found. Mitigated by deleting rather than disabling them, and by static test U9. \\
\bottomrule
\end{tabularx}
\end{table}

The current code is copied unchanged into the folder \code{pipeline/} of the project repository,
committed as a baseline, and
refactored on a branch \code{rerun-v2}. Unsafe paths are deleted, not switched off; the non-strict
mode ceases to exist.

\subsection{Module-by-module changes}

{\small
\begin{longtable}{@{}p{4.6cm} p{7.4cm} p{3.3cm}@{}}
\toprule
File & Change & Fixes \\
\midrule
\endfirsthead
\toprule
File & Change & Fixes \\
\midrule
\endhead
\bottomrule
\endfoot
\code{src/data/freeze_data.py} (new) & One-off download of M3/M4 into the cloud bucket with SHA-256 checksums; validation; M3 from \code{M3C.xls} with the Monash and Mcomp cross-checks; memory-safe M4 reader. & D2, m7 \\
\code{src/data/frozen.py} (new) & Canonical layout, content hash, verified loading. & D2 \\
\code{src/data/load_m_datasets.py} & Load only from the frozen copy; verify checksums; refuse a mismatch; one regular calendar from the position. & D2, F4, m6 \\
\code{src/data/eligibility.py} (new) & Eligibility table: length rule $L$, feature flags; refuses ineligible series downstream. & C4, M4, F1 \\
\code{src/data/validation.py} & \code{make_cutoffs} asserts eligibility instead of skipping short windows. & M4 \\
\code{src/features/compute_all.py} & Compute on $y_{1:n-3h}$ (line 66); a failed feature marks the series ineligible instead of writing default values. & C3, F1 \\
\code{src/features/nonlinear.py}, \code{arch.py} & Quarterly guards lowered to 20 (D5). & D5 \\
\code{src/features/mstl_features.py} & \code{decompose_series}: remove constant-trend and moving-average fallbacks; raise instead. & F2 \\
\code{src/data/sample_series.py} & 600 per source from the eligible pool; per-feature seed; strata recorded; 20\% minimum-support guard in bucket assignment. & M2, D7 \\
\code{src/models/baselines.py}, \code{component_wrappers.py} & Deleted. & C1 \\
\code{src/models/*_auto.py} & Split into \code{tune()} (Auto classes, Optuna, MASE) and \code{fit_predict()} (plain classes with a frozen config). Remove \code{_safe_lags}, \code{_safe_input_size} and all fallback branches. & C2, C4 \\
\code{src/models/statistical.py} & Approximate AutoARIMA kept; seasonal-naive fallback removed; per-series failures raise a recorded failure. & C1, m2 \\
\code{src/models/forecast_result.py} & Accept only \code{trained_*} sources and \code{failed} status. & C1 \\
\code{src/training/train_family.py} & Delete \code{_identity_scale}, \code{_apply_identity_adjustment} and \code{uid_not_in_global_result}; add component provenance, pairing keys, seed key and the statistical cache. & C1, m1, D17, D18 \\
\code{src/training/tune.py} & Real tuning: study creation, SQLite storage, 20-trial budget, frozen-config file with hash. & C2 \\
\code{src/training/strict_mode.py} & Removed; strictness is unconditional. & C1 \\
\code{src/training/task_graph.py} & Scopes \{cohort, low, medium, high\} plus seed-check tasks; delete \code{finetune_by_feature_bucket} and per-series modes. & M1, D17 \\
\code{src/training/run_experiment.py} & Split into stages \code{prepare}, \code{tune}, \code{evaluate}, \code{merge}, \code{gates}, \code{analyse}; hashes written at each stage; logs synced. & M1, m5 \\
\code{src/training/scheduler.py} & Shard selection by family $\times$ frequency; idempotent per-task outputs; rclone sync to the bucket. & D20 \\
\code{src/validation/gates.py} (new) & Gates G1--G13 (Section~\ref{sec:gates}). & m3 \\
\code{src/analysis/} (new) & Replaces \code{rebuild_paper_data.py} and the scratch scripts: loading with gates, deltas, Wilcoxon + Holm, MCM, Friedman + Nemenyi, D26 labels, tables, figures. & m3, m4 \\
\code{configs/rerun_v2.yaml} (new) & Encodes every decision in Section~\ref{sec:decisions}; the old budget modes that permit non-strict runs are removed. & --- \\
\end{longtable}
}

\subsection{Command-line interface}

\begin{verbatim}
python -m src.cli freeze-data --out <bucket>/data        # once, before R1
python -m src.cli prepare  --config configs/rerun_v2.yaml --run runs/v2
python -m src.cli tune     --run runs/v2 --shard <family>-<frequency>
python -m src.cli freeze   --run runs/v2          # writes configs_frozen.json + hash
python -m src.cli evaluate --run runs/v2 --shard <family>-<frequency>
python -m src.cli merge    --run runs/v2          # refuses mismatched hashes
python -m src.cli gates    --run runs/v2          # G1-G13, blocks analysis on failure
python -m src.cli analyse  --run runs/v2 --out <paper-repo>/results_v2
\end{verbatim}

%=====================================================================
\section{Testing and verification}\label{sec:testing}
%=====================================================================

Three layers: unit tests on code, integration tests on end-to-end runs, and gates on every
produced output. A failing gate blocks the next stage.

Tests inherited from the baseline that encode behaviour this protocol replaces are rewritten or
deleted in the commit that replaces that behaviour, with the reason in the commit message; every
other test passes at every commit. At the baseline four inherited tests fail: two expect the old
task graph (972 tasks; a per-bucket mode for statistical models) and two expect the old strict
leaderboard to reject zero or missing training times. They are resolved with the task graph and
provenance (D19) work.

\subsection{Unit tests}

{\small
\begin{longtable}{@{}p{0.7cm} p{10.8cm} p{3.6cm}@{}}
\toprule
ID & Asserts & Guards \\
\midrule
\endfirsthead
\toprule
ID & Asserts & Guards \\
\midrule
\endhead
\bottomrule
\endfoot
U1 & Perturbing any value after the first cutoff leaves all six features unchanged. & C3 \\
U2 & Series with fewer than $L$ pre-window-0 points are ineligible; \code{make_cutoffs} raises on them. & C4, M4 \\
U3 & On synthetic series of length exactly $L$, all six features return ``ok''. & D5, F1 \\
U4 & No sampled series carries a non-``ok'' feature flag. & F1 \\
U5 & STL: decomposition raises on too-short input; $T+S+R = y$ to $10^{-9}$; components unchanged when test data is perturbed. & F2, C3 \\
U6 & RelNaive, the cap and the seasonal naive match hand-computed examples. & metric \\
U7 & Sampling is deterministic given the seed, draws 600 per source, fills strata quotas, and uses different seeds per feature. & M2 \\
U8 & Tercile sizes differ by at most one; the 20\% support guard triggers on a degenerate synthetic feature. & D7 \\
U9 & Static check: no module can produce a non-trained forecast source; \code{baselines.py} is absent; the result object rejects other sources. & C1 \\
U10 & The same task run twice gives identical forecast hashes (CPU backends; neural with a fixed seed on CPU). & C1 jitter \\
U11 & Instantiated models carry exactly the frozen configuration, whatever the composition of the pool. & C4, C2 \\
U12 & Tuning data ends before each series' first cutoff; validation is the last $h$ points before it. & C2, C3 \\
U13 & STL-AC rows store three component provenances; any component failure fails the row. & m1 \\
U14 & Every tercile row has exactly one cohort twin; a missing twin fails the merge. & M1 \\
U15 & Merge refuses shards with a mismatched manifest, configuration, code, environment or data hash. & M1, m3 \\
U16 & Statistical forecasts are identical across feature samples for the same series and window. & D18 \\
U17 & A statistical per-series exception yields a failed row and the task continues; a global exception fails the task with no forecast rows. & D16 \\
U18 & Loading refuses data whose checksum differs from the frozen copy, a missing file, an unlisted source and a configuration that does not match the manifest; the loader contains no exception handler, downloader or random generator. & D2, F4 \\
U19 & Wilcoxon, Holm, Friedman, Nemenyi, Jonckheere--Terpstra and the D26 labelling match reference implementations and hand-computed toy cases. & D23, D26 \\
\end{longtable}
}

{\sloppy\textbf{Existing tests.} Kept: \code{test_rel_naive}, \code{test_component_recomposition},
\code{test_stl_no_leakage} (extended to U5), \code{test_checkpoint_resume},
\code{test_checkpoint_schema}, \code{test_feature_parallel_cache}, \code{test_timing_provenance}.
Updated: \code{test_representative_sampling}, \code{test_parallel_scheduler},
\code{test_global_model_adapters}, \code{test_feature_scaling_bins}, \code{test_benchmark_correctness}.
Rewritten: \code{test_feature_fallbacks} (fallback now means ineligible) and
\code{test_strict_provenance} (no fallback path exists). Replaced: \code{test_honest_optuna_status},
by U11 and U12.\par}

\subsection{Integration tests}

\begin{description}[leftmargin=1.1cm,style=sameline,font=\normalfont\bfseries]
  \item[I1] \textbf{Smoke:} synthetic data, all families, one trial per study, CPU. The full chain
  freeze-data $\rightarrow$ prepare $\rightarrow$ tune $\rightarrow$ freeze $\rightarrow$ evaluate
  $\rightarrow$ merge $\rightarrow$ gates $\rightarrow$ analyse passes.
  \item[I2] \textbf{Regression:} on the eligible subset of the 25 May strict run's manifest, the new
  code reproduces its AutoARIMA and AutoETS Direct metrics row for row ($|\text{diff}| < 10^{-6}$).
  \item[I3] \textbf{Analysis dry run:} the analysis package produces every table and figure of
  Table~\ref{tab:outputs} from pilot outputs.
  \item[I4] \textbf{Resume:} a task killed mid-run and resumed produces identical outputs.
  \item[I5] \textbf{API check (day 1):} every class and argument in Table~\ref{tab:models} and
  Section~\ref{sec:tuning} exists in the pinned environment.
\end{description}

\subsection{Output gates}\label{sec:gates}

{\small
\begin{longtable}{@{}p{0.7cm} p{10.4cm} p{4.0cm}@{}}
\toprule
ID & Passes if & Stage \\
\midrule
\endfirsthead
\toprule
ID & Passes if & Stage \\
\midrule
\endhead
\bottomrule
\endfoot
G1 & Every forecast row is \code{trained_*}; failures appear only in the failure table. & merge \\
G2 & All shards carry the frozen manifest hash. & merge \\
G3 & All shards carry the frozen configuration hash. & merge \\
G4 & All shards carry the same code commit (clean tree) and environment-lock hash. & merge \\
G5 & Task grid complete: expected $-$ completed = listed failures; global tasks have no unresolved failures; statistical per-series failures $\leq 0.5\%$, else investigated. & merge \\
G6 & Pairing complete for $\dstl$ (every STL row has a Direct twin) and $\dft$ (every tercile row has a cohort twin). & merge \\
G7 & Logged feature and tuning end indices all precede the first cutoff. & prepare, tune \\
G8 & No manifest series has fewer than $L$ pre-window-0 points. & prepare \\
G9 & No NaN or infinite metric; share of capped RelNaive $\leq 2\%$. & merge \\
G10 & STL-AC / STL-SN median training-time ratio in $[2, 4]$ per global family. & merge \\
G11 & Kolmogorov--Smirnov $p \geq 0.01$ between sample and eligible pool for each feature and M4 source; M3 reported. & prepare \\
G12 & Every populated bucket $\geq 20\%$ of its sample, or the feature is flagged degenerate. & prepare \\
G13 & Every shard carries the frozen data checksum. & prepare, merge \\
\end{longtable}
}

\subsection{Traceability}

\begin{table}[h]
\centering\small
\caption{Every defect is closed by a decision, a test and a gate.}\label{tab:trace}
\begin{tabular}{@{}lllll@{}}
\toprule
Defect & Decision & Unit tests & Integration & Gates \\
\midrule
C1 fallback stub, jitter & D16 & U9, U10, U17 & I1 & G1 \\
C2 no tuning & D12--D14 & U11, U12 & I5 & G3 \\
C3 feature leakage & D4 & U1, U5 & --- & G7 \\
C4 pool-dependent configs & D3, D14, D25 & U2, U11 & --- & G8 \\
M1 unpaired arms & D15, D20 & U14, U15 & I1 & G2, G6 \\
M2 no stratification & D6 & U7 & --- & G11 \\
M3 ``held-out series'' & D15 & --- & --- & paper text \\
M4 short histories & D3 & U2 & --- & G8 \\
M5 thin neural training & D14 & U11 & --- & --- \\
F1 feature fallback values & D4, D5 & U3, U4 & --- & G8 \\
F2 STL fallback & D9 & U5 & --- & --- \\
F3 code not versioned & D22 & --- & --- & G4 \\
F4 silent synthetic data & D2 & U18 & --- & G13 \\
m1 component provenance & D19 & U13 & --- & --- \\
m3 unchecked merge & D20 & U15 & I3 & G2--G6, G13 \\
m6 M4 time stamps & D2 & U18 & --- & --- \\
m7 M3 values & D2 & freeze cross-check & --- & G13 \\
\bottomrule
\end{tabular}
\end{table}

%=====================================================================
\section{Pilot}\label{sec:pilot}
%=====================================================================

\subsection{Scope}

\begin{itemize}
  \item Two features: evolving seasonality (the strongest earlier signal) and nonlinearity (it
  exercises the degenerate-bucket guard).
  \item Both frequencies; 100 series per source per feature (400 per feature).
  \item All 12 models, 3 strategies, 4 scopes, 3 windows, main seed only.
  \item Tuning: all 90 studies with 3 trials each, on a tuning set of 100 per source.
  \item One GPU machine, the same image and frozen data copy as the full run.
\end{itemize}

\subsection{What the pilot must show}

\begin{enumerate}
  \item All unit and integration tests pass, including I2 (regression) and I3 (analysis dry run).
  \item All gates G1--G13 pass on the pilot outputs, with G12 flagging nonlinearity as expected.
  \item Measured time per task and per trial, by family, frequency and strategy. These replace the
  estimates of Section~\ref{sec:runs} and set the shard allocation.
  \item Peak GPU memory per model at the largest configuration in each search space.
  \item Failure rate by family, with every failure explained.
\end{enumerate}

\subsection{Exit criteria (go / no-go)}

\textbf{Go} if all tests and gates pass, there is no unexplained failure, and the projected
wall-clock time of the full run on four machines is at most three days. The supervisor takes the
go / no-go decision (D30). The only permitted adjustments before going are:
\begin{itemize}
  \item narrowing a search space for memory or time;
  \item reducing trials from 20 to 15 for the slowest family (D12);
  \item re-balancing shards.
\end{itemize}
Each adjustment is recorded in the change log. Results from the pilot are never reported.

%=====================================================================
\section{Full runs}\label{sec:runs}
%=====================================================================

\subsection{Stages}

{\small
\begin{longtable}{@{}p{1.0cm} p{2.6cm} p{7.4cm} p{4.0cm}@{}}
\toprule
Stage & Name & Content & Output (frozen) \\
\midrule
\endfirsthead
\toprule
Stage & Name & Content & Output (frozen) \\
\midrule
\endhead
\bottomrule
\endfoot
R0 & Freeze code & Tag the commit (with the signed-off protocol); build and hash the environment image; data copy already frozen. & commit tag, image hash, data checksum \\
R1 & Prepare (CPU) & Eligibility on the full pool; features on pre-window-0 data for all eligible series (about 59,000); six samples; buckets; tuning set; cutoffs. Gates G7, G8, G11, G12, G13. & manifest bundle + hash \\
R2 & Tune & 90 studies, sharded (50 GPU, 40 CPU). Gate G7. & study files \\
R2b & Freeze configs & Best trial per study written and hashed. & \code{configs_frozen.json} + hash \\
R3 & Evaluate & Up to 1,488 tasks, sharded; statistical cache shared; cohort shards first, then tercile, then seed check. Each shard is validated on completion. & shard outputs + logs \\
R4 & Merge + gates & Merge; G1--G6, G9, G10, G13. Retry or fix any failed global task (D16). & result bundle + hash \\
R5 & Analyse & Analysis package on the frozen bundle. & tables, figures, reports \\
\end{longtable}
}

\subsection{Task grid}

\begin{itemize}
  \item \textbf{Cohort scope:} $6 \text{ features} \times 2 \text{ frequencies} \times 3
  \text{ strategies} \times 12 \text{ models} = 432$ tasks. Of these, 108 are statistical and are
  served by the cache.
  \item \textbf{Tercile scope:} $6 \times 2 \times 3 \times 9 \times 3 \text{ buckets} = 972$
  tasks. This falls to 918 if nonlinearity has an empty Medium bucket in both frequencies.
  \item \textbf{Seed check:} $1 \text{ feature} \times 2 \times 3 \times 7 \text{ stochastic models}
  \times 2 \text{ extra seeds} = 84$ tasks. Ridge and LinearRegression are deterministic and
  excluded.
  \item \textbf{Total:} up to 1,488 tasks, about 2.8 million forecast rows.
\end{itemize}

\subsection{Compute estimate and allocation}

The basis is the previous runs on one machine with 64 cores, one 24\,GB GPU and fixed small
configurations: 288 cohort-style tasks in about 12 hours, and 972 tercile tasks in about 50 hours.
At the new task count this is about 70 machine-hours. Tuned configurations may cost 1.5--3 times
more, giving roughly 105--210 machine-hours of evaluation. Tuning adds 50 GPU studies $\times$ 20
trials at an estimated 0.5--3 minutes each, i.e.\ 8--50 GPU-hours, plus CPU studies in parallel.

Initial allocation of the four machines (D20), re-balanced from the pilot's measured times:
\begin{itemize}
  \item machine A: neural monthly;
  \item machine B: neural and transformer quarterly;
  \item machine C: transformer monthly;
  \item machine D: ML and statistical, both frequencies (CPU-bound).
\end{itemize}
Expected wall-clock time: tuning 2--13 hours, evaluation 26--53 hours.

\subsection{Monitoring, failures and archival}

\begin{itemize}
  \item Heartbeat and status every minute; checkpoints every five minutes.
  \item Outputs and logs synced to the cloud bucket at every checkpoint via rclone, with checksums,
  so no machine holds the only copy (D20).
  \item A failed global task is retried once automatically. A second failure is investigated; the
  fix is committed, and the task alone is rerun with the same manifest and configuration (D16).
  \item \textbf{Priority:} cohort-scope shards (Experiment 1) run first, tercile-scope shards
  second and the seed check last, so an overrun cannot leave the main result incomplete.
  \item \textbf{Archival:} SHA-256 checksums of every file. The final bundle is read-only; every
  number in the paper is generated from it by committed scripts, and code, manifest and bundle are
  released with the paper (D29).
\end{itemize}

%=====================================================================
\section{Timeline}\label{sec:timeline}
%=====================================================================

Day 1 is the day of supervisor sign-off; the dates assume Tuesday 29 September 2026.

\begin{table}[h]
\centering\small
\caption{Two-week plan.}\label{tab:timeline}
\begin{tabularx}{\linewidth}{@{}l l X@{}}
\toprule
Day & Date & Work \\
\midrule
1 & Tue 29 Sep & Supervisor sign-off. Folder \code{pipeline/}; baseline commit; branch \code{rerun-v2}; environment lockfile; I5 API check; frozen data copy with checksums (U18). \\
2 & Wed 30 Sep & Eligibility, pre-window-0 features, sampling, buckets, manifest (U1--U8). \\
3 & Thu 1 Oct & Model wrappers split into tune / fit; stubs and jitter deleted; strict failures; provenance; seed keys (U9--U13, U16, U17). \\
4 & Fri 2 Oct & Task graph with scopes and seed check; sharding; merge; gates; analysis package incl.\ Friedman / Nemenyi and D26 labels, tested on synthetic data (U14, U15, U19, I1, I4). \\
5 & Sat 3 Oct & Pilot on one GPU; I2, I3; gates; cost projection; go / no-go by the supervisor. \\
6 & Sun 4 Oct & R0; R1 (prepare); start R2 (tuning) on four machines. \\
7 & Mon 5 Oct & Finish R2; R2b freeze; start R3 (evaluation), cohort shards first. \\
8--9 & Tue 6 -- Wed 7 Oct & R3; validate shards as they finish. \\
10 & Thu 8 Oct & R4 merge and gates; fix and rerun any failed task; buffer. \\
11 & Fri 9 Oct & R5 analysis; tables and figures. \\
12--14 & Sat 10 -- Mon 12 Oct & Rewrite methods and results; supervisor review. \\
\bottomrule
\end{tabularx}
\end{table}

The critical path is days 1--5. A one-day slip in implementation removes one day of writing. The
pilot's go / no-go falls on a Saturday; if the supervisor is not available that day, R0 and R1 may
start, since they do not depend on the decision, but tuning waits for it. A no-go costs one to two
days and is escalated to the supervisor the same day.

%=====================================================================
\section{Risks}\label{sec:risks}
%=====================================================================

{\small
\begin{longtable}{@{}p{4.6cm} p{2.0cm} p{8.6cm}@{}}
\toprule
Risk & Likelihood & Mitigation \\
\midrule
\endfirsthead
\toprule
Risk & Likelihood & Mitigation \\
\midrule
\endhead
\bottomrule
\endfoot
Auto-class APIs differ in the pinned versions & medium & I5 on day 1. Fallback: an Optuna loop around the plain classes with the same spaces and objective, recorded in the change log. \\
Tuning overruns & medium & Measured in the pilot; reduce trials to 15 for the slowest family before R2 (D12). \\
Evaluation overruns & medium & Cohort shards first, seed check last; re-balance machines. \\
Machine preemption & medium & Per-task checkpoints; idempotent tasks; continuous sync to the bucket. \\
GPU non-determinism & high (small) & Fixed seeds, deterministic flags, forecast hashes; quantified by A15. \\
Degenerate buckets beyond nonlinearity & low & G12 flags them; handled as in D7. \\
Residual leakage & low & U1, U5, U12, G7. \\
Statistical failures on short series & low & Eligibility should remove them; U17 records any that remain. \\
Supervisor unavailable at a gate & low & R0--R1 may proceed; tuning and evaluation wait for sign-off. \\
Results contradict earlier drafts & high & Expected; the analysis plan is frozen before the run, so the conclusions follow the data. \\
\end{longtable}
}

%=====================================================================
\section{Consequences for the paper}\label{sec:paper}
%=====================================================================

\begin{itemize}
  \item \textbf{Methods to rewrite:} frozen data (D2), eligibility (D3), pre-window-0 features
  (D4, D5), per-feature sampling with M3 close to a census (D6), buckets and the support guard (D7),
  real tuning with its spaces and budget (D12--D14, D25), the two training scopes as a temporal
  hold-out on the same series (D15), strict execution and failure accounting (D16).
  \item \textbf{Sample sizes change:} 2,400 series per feature (1,200 per frequency, about 400 per
  bucket), about 8,500 distinct series overall.
  \item \textbf{All results sections are regenerated} from the frozen bundle. No number from any
  earlier version is carried over.
  \item \textbf{Language of effects:} ``gain'', ``loss'', ``negligible'' and ``no evidence'' are used
  exactly as defined in Section~\ref{sec:practical}.
  \item \textbf{New figures:} critical-difference diagrams for both scopes (H7).
  \item \textbf{Cross-feature statements} (``feature A moderates more than feature B'') compare
  different samples and are worded as such; each feature's gradient is reported on its own sample.
  \item \textbf{Limitations:} one main seed (quantified by A15); STL only; four features computed
  from STL components; M3 Quarterly close to a census; the lowered quarterly feature guards (D5,
  with A7); no feature-based selector (D28).
  \item \textbf{Reproducibility statement:} code, manifest and frozen results released (D29).
  \item \textbf{Unchanged:} introduction, related work, references, problem formulation, and the
  documented failure mode of tercile bucketing (now detected by G12).
\end{itemize}

\appendix

%=====================================================================
\section{Defects found in the review}\label{app:defects}
%=====================================================================

{\small
\begin{longtable}{@{}p{0.8cm} p{8.8cm} p{5.6cm}@{}}
\toprule
ID & Defect & Evidence \\
\midrule
\endfirsthead
\toprule
ID & Defect & Evidence \\
\midrule
\endhead
\bottomrule
\endfoot
C1 & Global baselines and statistical STL results are stub forecasts from a non-strict run; forecasts jittered per task. & \code{baselines.py:30-38}; \code{train_family.py:53-70}; run summary ``NOT BENCHMARK-VALID''. \\
C2 & No tuning; fixed configurations; placeholder Optuna files. & \code{mlforecast_auto.py:28-43}; \code{neural_auto.py:34}; \code{tune.py}; empty \code{optuna/}. \\
C3 & Features computed on $y_{1:n-h}$, seeing test windows 0--1. & \code{compute_all.py:66}. \\
C4 & Seasonal lag and input size set by the shortest series in the pool. & \code{mlforecast_auto.py:45-50}; \code{neural_auto.py:27-29}; 33\% of quarterly cohort fits had the lag. \\
M1 & Trained global and tercile results on different series; no global STL-AC. & 1,681 shared series; run configuration of 25 May. \\
M2 & Candidate pool equal to the quota, so no stratification took place. & 23 May configuration: pool cap 750. \\
M3 & ``Held-out series'' wording; evaluation is a temporal hold-out. & \code{train_family.py:192-206}. \\
M4 & Training histories as short as 6 quarterly and 12 monthly points. & \code{validation.py:6-16}. \\
M5 & Thin neural training (500 steps, no explicit scaling). & \code{neural_auto.py:32-48}. \\
F1 & Feature fallback values (0.0 / 1.0) assigned to short series. & 1.7--2.9\% of the 23 May sample. \\
F2 & STL silently replaced by a constant trend or a moving average. & \code{mstl_features.py:28-51}. \\
F3 & Pipeline source not under version control. & 0 tracked files in \code{forecasting_pipeline/src}. \\
F4 & The data loader caught every exception and substituted synthetic series (for M4 first a local copy of the training file only, without the test period), without recording it; the smoke budget forced synthetic data. Latent: a scan of every stored run found synthetic series only in the six smoke runs (160 each); all pilot and large runs and all series behind the paper's results (10,659 global-arm and 3,000 bucket-arm series) are real M3/M4 series. Found while implementing D2 (v1.3). & \code{load_m_datasets.py:50-85} and \code{run_experiment.py:66-67} at the baseline commit; scan of 14 feature tables and of \code{paper_workspace/final_results}. \\
m1 & STL-AC provenance taken from the trend model only. & \code{train_family.py:264-271}. \\
m2 & Statistical known-limit fallback even in strict mode. & \code{statistical.py:64-83}. \\
m3 & Rebuild merges runs without provenance checks. & \code{rebuild_paper_data.py:31}. \\
m4 & Analysis scripts not versioned. & --- \\
m5 & No logs kept for the 28--30 May runs. & --- \\
m6 & M4 time stamps (integer positions) were converted to nanosecond date-times, so a task's inferred frequency depended on which series came first. No effect on forecasts found (global training frames are rebuilt on a regular calendar and no model uses calendar inputs), but fragile. Found while implementing D2 (v1.3). & \code{load_m_datasets.py:88}; \code{train_family.py:235}; \code{global_adapter_data.py:26}. \\
m7 & M3 values of the earlier runs were rounded to single precision by the TSF reader of \code{datasetsforecast} (29,051 monthly and 14,983 quarterly values; relative error ${\le}\,5.7\times10^{-8}$, i.e.\ the eighth significant digit). Found while implementing D2 (v1.3). Corrected in v1.4: v1.3 also listed the value $-1200$ of N2786 at $t=84$ as a sign error in the Monash file; it is the value of the original \code{M3C.xls}. & Every value compared with \code{M3C.xls}; \code{np.float32} in \code{datasetsforecast/utils.py}. \\
\end{longtable}
}

%=====================================================================
\section{Decision record: alternatives considered}\label{app:alternatives}
%=====================================================================

The questions put to D.\ Costa on 28 September 2026 (rows marked v1.3 or v1.4: 29 September), the option chosen, and the alternatives
rejected. A rejected alternative may only be adopted later through the change log.

{\small
\begin{longtable}{@{}p{0.8cm} p{3.2cm} p{4.0cm} p{7.0cm}@{}}
\toprule
Dec. & Question & Chosen & Rejected alternatives (and why they were on the table) \\
\midrule
\endfirsthead
\toprule
Dec. & Question & Chosen & Rejected alternatives (and why they were on the table) \\
\midrule
\endhead
\bottomrule
\endfoot
D1 & Reuse old results? & Full rerun & Reuse the 25 May strict run and add missing arms (different samples, C3/C4 remain). \\
D2 & Datasets & M3+M4 monthly and quarterly & Add M4 yearly as a negative control; add weekly/daily (needs MSTL). \\
D2 & Data source & Frozen copy with checksums & Download at run time (a changed mirror would change the data). \\
D2 & M3 source (v1.4) & Original \code{M3C.xls} (IIF); Monash and Mcomp as cross-checks & Monash as the source (identical values, but a redistribution); Mcomp~2.7 (differs from the original in one value without documentation; chosen in v1.3 on the mistaken premise that Monash had an error); Monash through \code{datasetsforecast}, as in the earlier runs (single-precision rounding). \\
D2 & M3 ids (v1.3) & Official M3 ids (\code{N1402}, \dots) & Positional ids of the earlier runs (\code{M1}, \dots), which do not match the literature. \\
D5 & Short quarterly series & A: guards lowered to 20 & B: two quarterly windows; C: original guards (M3 Quarterly $\approx$116); D: 400 per source. \\
D4 & When features are computed & Once, before window 0 & Recomputed every window (bucket membership would change). \\
D6 & Sample pool & All eligible series & Random 5,000 per M4 source (faster). \\
D6 & Sampling design & Per feature, 600 per source & One common 3,000-series sample (pairable across features). \\
D7 & Nonlinearity feature & Keep definition, flag it & Make it continuous; break ties at random. \\
D7 & Degeneracy threshold & Any bucket $< 20\%$ & $< 10\%$; only an empty bucket. \\
D25 & ML inputs & Lags $1\dots m$ & Lags + calendar dummies; lags + rolling statistics. \\
D10 & Model set & Keep the 12 & Add DLinear; add DLinear + a zero-shot foundation model. \\
D11 & ARIMA search & Approximate & Exact search (slower; no regression reference). \\
D12 & Tuning approach & Real light tuning (Auto classes) & Fixed configurations; tune once per family on Direct only. \\
D12 & Tuning objective & MASE & RelNaive on validation (custom loss); MAE (not scale-free). \\
D12 & Trials & 20; 15 if pilot overruns & 20 fixed; 30 fixed. \\
D13 & Where to tune & Once, on pre-test data & Separately per feature sample. \\
D13 & Validation data & One block of $h$ & Two blocks monthly, one quarterly. \\
D13 & STL-AC tuning & Separate study per component & One shared configuration; reuse the Direct configuration. \\
D15 & Bucket-model tuning & Reuse the global configuration & Re-tune per bucket (confounds pool size with tuning). \\
D15 & Experiment 2 scope & Fully paired & Reduced scope; move to the follow-up paper. \\
D15 & Training pool & M3 + M4 pooled & Separate models per source. \\
D17 & Seeds & One seed + seed check & One seed only; three seeds for neural and transformer. \\
D8 & Refitting & Every window & Fit once, roll forward. \\
D16 & Persistent failures & Fix and rerun the task & Exclude the configuration and report it. \\
D20 & Machines & Four & Two; six or more. \\
D20 & Sync target & Cloud bucket & University server; Google Drive. \\
D22 & Library versions & Pin the May versions & Upgrade to the latest. \\
D22 & Repository & New folder \code{pipeline/} in the existing project repository, branch \code{rerun-v2} (amended in v1.2) & A new dedicated repository (the v1.1 choice, replaced on instruction of D.\ Costa). \\
D29 & Release & Code, manifest and results & Code only; nothing until the follow-up. \\
D21 & Primary metric & $\RN$ & MASE; both co-primary. \\
D26 & Practical threshold & 0.01 & 0.02; no threshold. \\
D23 & Inference & Wilcoxon + Holm + MCM + Friedman/Nemenyi & Wilcoxon only; Bayesian signed-rank with ROPE. \\
D23 & CD-diagram scope & Cohort and tercile separately & Cohort only; families instead of configurations. \\
D27 & Family pooling & Equal weights & Weight by model count; never pool. \\
D27 & Descriptive level & Instances; tests per series & Series level everywhere. \\
D28 & Selector & None in this paper & A simple exploratory rule. \\
D24 & Freeze & Tagged commit & Public pre-registration on OSF. \\
D30 & Sign-off & Supervisor & Supervisor + partner; D.\ Costa alone. \\
\end{longtable}
}

%=====================================================================
\section{Change log}\label{app:changelog}
%=====================================================================

\begin{longtable}{@{}p{1.3cm} p{1.9cm} p{11.8cm}@{}}
\toprule
Version & Date & Change \\
\midrule
\endfirsthead
\toprule
Version & Date & Change \\
\midrule
\endhead
\bottomrule
\endfoot
1.0 & 28 Sep 2026 & First draft for approval. \\
1.1 & 28 Sep 2026 & All open decisions taken in eight rounds of questions: D5 confirmed (option A); new decisions D25--D30 (ML inputs, practical threshold, aggregation, selector, release, sign-off); D2, D8, D11--D13, D15--D17, D20--D24 made explicit; Friedman + Nemenyi rankings (H7) and seed check (A15) added; gates G13 and tests U18--U19 added; timeline moved to start on the day of sign-off; Appendix~\ref{app:alternatives} added. \\
1.2 & 28 Sep 2026 & D22 amended on instruction of D.\ Costa: the pipeline lives in the new folder \code{pipeline/} of the existing project repository (\code{github.com/DiogoCosta18/TimeSeriesForecasting}), branch \code{rerun-v2}, instead of a new dedicated repository. Consequence: sharing the code with the partner university means sharing the repository, paper workspace included. Sections 1.2, 4.6, 6.1, the timeline and Appendix~\ref{app:alternatives} updated accordingly. \\
1.3 & 29 Sep 2026 & D2 amended on instruction of D.\ Costa after the M3 comparison: M3 from Mcomp~2.7 (CRAN) with official ids, the Monash copy kept as a bit-for-bit cross-check; rdata 0.11.2 and xarray 2026.7.0 added to the lock (no other pin changed); Appendix~\ref{app:alternatives} extended. Recorded while implementing D2: defects F4 (silent synthetic-data fallback in the loader; shown to be latent), m6 (M4 time stamps) and m7 (M3 values of the earlier runs) added to Appendix~\ref{app:defects} and to the traceability table; Section~\ref{sec:data} describes the frozen copy as implemented (sources, validation, file layout, manifest, order-only calendar); U18 extended; the rule for inherited tests and the four failing baseline tests stated in Section~\ref{sec:testing}; module table updated. \\
1.4 & 29 Sep 2026 & Correction of 1.3. The statement that the Monash file carries a sign error was wrong: the original competition file \code{M3C.xls} (IIF), obtained the same day, has $-1200$ at N2786, $t=84$, and Monash reproduces it exactly; Mcomp~2.7 is the copy that differs ($+1200$, undocumented). D2 decided again by D.\ Costa: M3 from \code{M3C.xls} (pinned SHA-256), with Monash (identical) and Mcomp (one documented difference) as cross-checks that stop the freeze on any other difference. Section~\ref{sec:data}, the decision register, defect m7, the environment paragraph, the module table and Appendix~\ref{app:alternatives} corrected accordingly. \\
1.5 & 1 Oct 2026 & \textbf{Sign-off:} v1.4 approved by the supervisor (reported by D.\ Costa on 1~October 2026); the approved text is the commit tagged \code{protocol-v1.4-approved}. From this version on, the body keeps the approved text and amendments are listed here only. \textbf{Amendments:} (a)~D20: the cloud bucket is a private Backblaze B2 bucket, accessed with rclone (decision of D.\ Costa, 1~October); credentials are created by D.\ Costa and never stored in the repository. (b)~Section~6.3: the frozen data copy is created and checked with \code{python -m src.data.freeze_data freeze -{}-out DIR -{}-code-commit SHA -{}-m3c M3C.xls} and \code{... verify -{}-data-dir DIR} (already done: copy \code{m3m4-v1}), then uploaded to the bucket with rclone; the \code{freeze-data} line of the listing becomes a \code{src.cli} subcommand with the same arguments when the stages are split (day~4). (c)~Section~\ref{sec:env}: the lock also pins sympy 1.13.1, the version torch 2.5.1 requires (the 21 May freeze had 1.14.0, which came with a later torch); \code{pipeline/environment/README.md} gives the derivation. Appendix~\ref{app:changelog} is typeset as a table that can break across pages. \\
1.6 & 3 Oct 2026 & Implementation of days 1--2; no decision changed. (a)~D20: the frozen copy is in the private Backblaze B2 bucket (path \mbox{\code{data/frozen_m3m4_v1/}}, 235\,MB), verified by rclone (SHA-1, then byte for byte) and by \code{freeze_data verify} on a fresh download. (b)~D4, F1, F2 as implemented: besides the fallbacks named in Section~6.2, the feature code substituted values on short or constant input, computed other statistics when a library call failed (moments, squared autocorrelation, and a full scan for ruptures that kept the flag ``ok''), interpolated missing values, set seasonal strength to 0 when undefined, and turned series exceeding a 120\,s timeout into fallbacks, so that eligibility depended on machine speed. All removed: a feature is ``ok'', ``undefined:\emph{reason}'' or ``error:\emph{type}'', with no value otherwise, and there is no time limit. On real data every ``ok'' value equals the earlier code bit for bit except the D5 cases. Consequence of F2: a series with a constant rolling window (36 monthly, 16 quarterly points) has no evolving-seasonality value and is ineligible. A failed parquet write no longer produces a CSV file under a parquet name. (c)~Seeds: ``seed $= s_0 + k$'' ($s_0 = 20260521$) is the numpy seed tuple $(s_0 + k, m, \text{source})$, one stream per feature, frequency and source; the tuning set uses $(s_0, 1000, m, \text{source})$. (d)~Buckets: the cut points are the values at positions round$(n/3)$ and round$(2n/3)$ (ties by series id), tied values share a bucket, and an empty bucket counts as below 20\% (Appendix~\ref{app:alternatives}: ``any bucket $< 20\%$''). A pool below 600 stops the stage; nothing is filled from elsewhere. (e)~Verification run of stage R1 at commit \code{cb5c047} (bundle \code{de54ccf5}\dots): eligible M3 Monthly 1,039, M4 Monthly 35,010, M3 Quarterly 634, M4 Quarterly 22,591, in total 59,274 (Table~\ref{tab:eligible} estimated 1,040, 35,000, 634, 22,500). Ineligible after the length rule: 110 M4 Monthly series (one constant throughout, 109 through evolving seasonality only) and 21 M4 Quarterly. 8,492 distinct sampled series; nonlinearity degenerate in both frequencies (Medium empty; Low 1,007 monthly and 976 quarterly), as Section~3.5 expects; every other feature 400/400/400; Kolmogorov--Smirnov $p \geq 0.9999$ for every sample (G11). (f)~Until the stage CLI of Section~6.3 exists, the stage runs as \code{python -m src.stages.prepare}. \\
1.7 & 3 Oct 2026 & Implementation of day 3; no decision changed. (a)~The model layer is the new package \code{src/forecast/} (registry, the search spaces of Table~\ref{tab:spaces}, result object, fits from frozen configurations, tuning, metrics, evaluation engine). The legacy modules of Section~6.2 (\code{src/models/*}, \mbox{\code{train_family.py}}, \code{baselines.py}, the component stub, the jitter, strict mode) are deleted when the runner switches to it on day~4. (b)~Global models receive the positions $t$ as time stamps (frequency 1). (c)~Tuning: the neural training loss is MAE, the models' default (the validation loss is MASE as specified); early stopping uses a patience of 5 validation checks every 50 steps, values the protocol does not fix; the frozen configuration holds the number of steps the best configuration actually trained, read from its refit on the same validation block; the TPE sampler is seeded per study from $s_0$ and every fit uses $s_0$ (D17); for the STL targets, STL is fitted once on each tuning series up to its first cutoff; every trial is archived in SQLite with its parameters, validation MASE, timestamps and duration. (d)~Determinism (U10): forecasts are reproducible bit for bit on CPU. RandomForest runs on one thread (with several, the tree predictions are summed in completion order and the last bits vary) and XGBoost on a fixed four. On a GPU there is no deterministic CUDA median (robust scaler), so the deterministic flag is ``warn'' there and the seed check (A15) quantifies the variation; the device is recorded with every forecast. (e)~Statistical models are fitted with fit and predict, which give the earlier runs' \code{forecast} results bit for bit, so regression test I2 is unaffected. (f)~POCID is computed as the paper's equation. Defect m8: the code that produced the paper's POCID values (\code{rebuild_paper_data.py}) used $h-1$ steps without the anchor on the last training value and counted a zero change on both sides as agreement, so the published values do not follow the published equation. (g)~Smoke run on real data (quarterly, one feature sample of 1,200 series, laptop RTX~3050~Ti): no failures; Ridge 1.5--9\,s per task of three windows, NHITS on GPU 31\,s, statistical fits 0.03--0.18\,s per series and window. \\
\end{longtable}

\end{document}
